\documentclass[11pt]{article}
\usepackage[margin=2.1cm]{geometry}
\usepackage{booktabs}
\usepackage{amsmath}
\usepackage{enumitem}
\usepackage[T1]{fontenc}
\usepackage{lmodern}

\setlength{\parindent}{0pt}
\setlength{\parskip}{5pt}
\newcommand{\code}[1]{\texttt{#1}}
\newcommand{\answerbox}[1]{\vspace{2pt}\fbox{\parbox[c][#1][c]{\dimexpr\linewidth-2\fboxsep-2\fboxrule}{\ }}\vspace{4pt}}

\begin{document}

\begin{center}
{\Large\bfseries EE\,362 --- Pre-Lab 12}\\[2pt]
{\large CPU Integration}\\[6pt]
\begin{tabular}{@{}p{6.5cm}p{4cm}p{4cm}@{}}
Name: \hrulefill & Section: Mon / Tue & Date: \hrulefill
\end{tabular}
\end{center}

\vspace{2pt}\hrule\vspace{4pt}

{\small\textbf{Due at the start of the lab session.} Five points. \textbf{Bring
your printed control table} --- you will need it during the lab, and Question 2
uses it.}

\vspace{4pt}\hrule

\section*{Question 1 --- Hand-compute the trace}

Fill in the first six cycles of the program by hand, from the ROM listing and
your control table. Do not simulate first --- the point is that you can predict
the machine.

\begin{center}
\small
\begin{tabular}{clllll}
\toprule
cyc & PC & instruction & ALU result & register written & next PC \\
\midrule
0 & \code{0x00} & \code{LUI r1, 11} & \rule{2.2cm}{0.4pt} & \rule{2.4cm}{0.4pt} & \rule{1.6cm}{0.4pt} \\[6pt]
1 & \code{0x04} & \rule{2.6cm}{0.4pt} & \rule{2.2cm}{0.4pt} & \rule{2.4cm}{0.4pt} & \rule{1.6cm}{0.4pt} \\[6pt]
2 & \code{0x08} & \rule{2.6cm}{0.4pt} & \rule{2.2cm}{0.4pt} & \rule{2.4cm}{0.4pt} & \rule{1.6cm}{0.4pt} \\[6pt]
3 & \code{0x0C} & \rule{2.6cm}{0.4pt} & \rule{2.2cm}{0.4pt} & \rule{2.4cm}{0.4pt} & \rule{1.6cm}{0.4pt} \\[6pt]
4 & \code{0x10} & \rule{2.6cm}{0.4pt} & \rule{2.2cm}{0.4pt} & \rule{2.4cm}{0.4pt} & \rule{1.6cm}{0.4pt} \\[6pt]
5 & \code{0x14} & \rule{2.6cm}{0.4pt} & \rule{2.2cm}{0.4pt} & \rule{2.4cm}{0.4pt} & \rule{1.6cm}{0.4pt} \\
\bottomrule
\end{tabular}
\end{center}

\textbf{(b)} The loop body is four instructions and runs a fixed number of
times, plus a partial final pass. Work out the total instruction count.

\vspace{4pt}
loop passes $=$ \rule{2cm}{0.4pt} \hfill total instructions $=$ \rule{2cm}{0.4pt}
\hfill CPI $=$ \rule{1.5cm}{0.4pt}
\vspace{6pt}

\answerbox{1.6cm}

\section*{Question 2 --- The register file connections}

The register file has two read ports. One is wired to the \code{rs} field and
one to \code{rt}. Getting them the wrong way round still lets \emph{three} of
the four instruction types work.

\textbf{(a)} For each instruction, say which register must arrive on
\code{Aout} and which on \code{Bout}, and why.

\begin{center}
\begin{tabular}{lp{4.2cm}p{6.4cm}}
\toprule
instruction & \code{Aout} / \code{Bout} & why \\
\midrule
\code{ADDI rt, rs, imm} & \rule{4cm}{0.4pt} & \rule{6.2cm}{0.4pt} \\[7pt]
\code{SLL rd, rt, sa} & \rule{4cm}{0.4pt} & \rule{6.2cm}{0.4pt} \\[7pt]
\code{SW rt, off(rs)} & \rule{4cm}{0.4pt} & \rule{6.2cm}{0.4pt} \\
\bottomrule
\end{tabular}
\end{center}

\textbf{(b)} Conclude: \code{Asel} $=$ \rule{2.5cm}{0.4pt} \quad and \quad
\code{Bsel} $=$ \rule{2.5cm}{0.4pt}

\textbf{(c)} If you swapped them, which single instruction in this program would
break first, and what would the symptom be in the trace?

\answerbox{2.0cm}

\section*{Question 3 --- The debugging method}

Your CPU runs and the final stored word is wrong.

\textbf{(a)} Write down the four checks, in order, and say what each one tells
you if it fails.

\answerbox{4.2cm}

\textbf{(b)} Suppose checks 1, 2 and 3 all pass. What have you learned, and
where must the fault be?

\answerbox{1.8cm}

\textbf{(c)} Why is it worth doing the checks in order even when you are already
fairly sure what is wrong?

\answerbox{1.6cm}

\vfill
{\small\emph{Question 1 is the one that matters this week.} You cannot debug a
machine you cannot predict --- and this lab is entirely debugging.}

\end{document}
